\section{Odd nonzero helper quotients and cofinal signed-cap
compositions}\label{odd-nonzero-helper-quotients-and-cofinal-signed-cap-compositions}

\subsection{Statement, exact formula, and
attribution}\label{statement-exact-formula-and-attribution}

For each published endpoint let R\_j be exactly its residual list in
\texttt{equation.json}, with the declared integer-coordinate order. The
equation is

\begin{verbatim}
P(z_0,...,z_(n-1)) = sum_j R_j(z)^2 = 0, all z_i in Z.
\end{verbatim}

\texttt{expanded\_polynomial.json} is its complete collected coefficient
list, not a circuit estimate or a polynomial with variable exponents.
The degree is four. The proof below establishes zero existence iff a
finite ZFC refutation exists, at the inherited ordinary C222/refutation
calibration. Thus ZFC proves that the no-integer-zero sentence is
equivalent to Con(ZFC). This is not a proof of consistency, a new
kernel-checked ZFC derivation, or a global minimum claim.

The starting body is the complete fused-nine signed-cap source at
upstream commit 61f87277eb8acdd9140021a346b7a869d9cbb32a, with 406 body
coordinates, 341 body rows, 475 guarded coordinates and 364 guarded
rows. Its full literal files are retained in \texttt{source/}; all
copied file hashes are in SOURCE\_PINS.

The contribution is the odd-helper quotient below, its signed-base
proof, and its actual composition with the accepted
signed-target/fused-factor source. The two support products are upstream
623d491b27c719151f652f15269bf65d10065db7. The odd frame, cofinal square
horizon, consecutive-product beta modulus and self-sized radix are
credited to incoming source131 and its cofinal sources. The old round16
odd cap already used an odd auxiliary quotient in a different
positive-target construction. No historical-priority claim is made for
the underlying parity trick. Exact new materialized compositions, not
subtraction of advertised savings, determine the endpoint counts.

\subsection{1. Literal source changes and
conventions}\label{literal-source-changes-and-conventions}

All coordinate numbers here refer to the unprojected 406-coordinate
body. Use t=v0+16, x=v1, y=v2, F=v3, B=v5, V0=v6, r=v7, s=v8, r0=v11,
N=v26, q=v179, W=v245, A=v246, z=v247. The helper is H=v250, J=v251,
d=v252, f=v253, AJ=v254, Dminus1=v255, Hf=v256, e=v375. The cap is
U=v386, C=v387, AC=v388, C2=v389, S=v390, R=v391, DR=v392, g=v393,
h=v394.

Replace only these multiplier terms in the body:

\begin{verbatim}
v248=W(z+1)       -> v248=W(2z+1),
J=C2(e+1)         -> J=C2(2e+1).
\end{verbatim}

Keep all other frame, cap, primary norm, pool and power rows. Delete
only the frame guard z\textgreater=0 and helper guard
Dminus1-4\textgreater=0, including their private triples. Replace t0's
ternary graph by t0=h0\^{}2. Replace the scalar guard pairs (183,184)
and (191,192) by graphs for their products. The \texttt{cofinal\_y}
sibling also adds y=V0(V0-1), replaces (213,214) by a product graph, and
changes the strict fused graph (r0+1)P8-1\textgreater=0 to
P8-1\textgreater=0. Its old row y=F+r0+1 remains.

The optional \texttt{self\_radix} sibling is described in section 7.
Every other change is a recorded private-witness reindexing, Horn
identity, integral monic elimination, or integer translation. The
transformation certificate lists every changed full body row and every
retained/deleted/fresh graph.

For any integer c, c(c+1)\textgreater=0. Legendre's three-square theorem
applied to 4n+1 shows that n\textgreater=0 iff n=a\^{}2+b\^{}2+c(c+1)
for some integers a,b,c. Each new graph is quadratic and has its own
private witnesses.

\subsection{2. The odd frame requires no sign
predicate}\label{the-odd-frame-requires-no-sign-predicate}

The frame remains W=2(2+sum of squares of all distinct actual power
atoms), so W is even, W\textgreater=516, every atom has absolute value
below W, and there are 15 primary indices with absolute values below W.
The t target is nonzero, because t=16+h0\^{}2. The unprojected scale
rows are exactly

\begin{verbatim}
T=W(2z+1), E=(W+1)T, A^2+T^2-E^2=1.
\end{verbatim}

Thus A\^{}2-((W+1)\^{}2-1)T\^{}2=1 and T is nonzero. Write
a=\textbar A\textbar. Pell classification at base W+1 gives
\textbar T\textbar=Y\_n(W+1) for n\textgreater=1. Reduction modulo W
gives Y\_n(W+1)=n mod W, so W\textbar n.~Consequently

\begin{verbatim}
a>=X_W(W+1)>W^W,  a-1>4W^3.
\end{verbatim}

These statements use no requested exponential meaning or unproved sign.

For completeness choose A=X\_W(W+1). Because W+1 is odd, X\_n(W+1) is
odd for all n, Y\_n(W+1) is odd when n is odd, and
Y\_(2n)=2X\_nY\_n.~Induction on the power of two in n gives
v2(Y\_n(W+1))=v2(n). Together with W\textbar Y\_W(W+1), the quotient
Y\_W(W+1)/W is a positive odd integer. Set z=(Y\_W/W-1)/2. This supplies
the odd frame for every fixed semantic front, not only samples.

\subsection{3. Odd helper theorem, with signed A, C and comparison
base}\label{odd-helper-theorem-with-signed-a-c-and-comparison-base}

Retain the three norms and congruences

\begin{verbatim}
U^2-(A^2-1)C^2=1, H^2-(A^2-1)J^2=1,
S^2-(D^2-1)R^2=1,
J=C^2(2e+1), D=1+4J(d+1)=A+Hf,
S=U+Hg, R=W+4Ch.
\end{verbatim}

Also retain the quadratic cap and primary-only pooled congruence

\begin{verbatim}
C^2-1-sum_i y_i^2 = a0^2+b0^2+c0(c0+1),
sum_i(y_i-k_i)^2=(A-1)Q.
\end{verbatim}

The first gives C nonzero. If \textbar C\textbar=1, all y\_i=0, and the
pool would make the strictly positive sum sum k\_i\^{}2, smaller than
W\^{}3, divisible by A-1, whose absolute value exceeds 4W\^{}3. This is
impossible. Hence L=\textbar C\textbar\textgreater1. Classification
gives L=Y\_i(a), \textbar U\textbar=X\_i(a), i\textgreater=2. In
particular L\textgreater=2a\textgreater W and 0\textless i\textless L.
This excludes the exceptional index-one endpoint explicitly.

Since 2e+1 is nonzero,
\textbar J\textbar\textgreater=C\^{}2=L\^{}2\textgreater1.
Classification of the helper gives \textbar J\textbar=Y\_l(a),
\textbar H\textbar=X\_l(a), l\textgreater=2. In particular
\textbar H\textbar\textgreater=2a\^{}2-1\textgreater a+1\textgreater=\textbar A-1\textbar.
If d+1=0, then D=1 and H divides A-1, contradicting this strict
inequality. Therefore d+1 is nonzero and

\begin{verbatim}
|D|=|1+4J(d+1)|>=4|J|-1>=15.
\end{verbatim}

No sign of J, d, D, A or C has been inferred or assumed.

The inherited Pell squared-divisibility fact, applied to L\^{}2 dividing
\textbar J\textbar, yields i\textless=l and L\textbar l. The congruence
R=W modulo4C, with0\textless W\textless L, implies R is nonzero.
Classify the comparison norm at \textbar D\textbar, with absolute index
j\textgreater=1. For negative D the polynomial recurrences satisfy
X\_j(D)=sgn(D)\^{}j X\_j(\textbar D\textbar) and
Y\_j(D)=sgn(D)\^{}(j-1)Y\_j(\textbar D\textbar); the independent
coordinate signs absorb these factors. Thus the exact polynomial
congruence D=A mod H and S=U mod H gives X\_j(a)=+/-X\_i(a) mod X\_l(a),
hence j=+/-i mod 2l by the same inherited up-to-sign residue lemma. The
conditions i\textgreater=2 and i\textless=l are already established.

Since D=1 mod 4L and R=W mod 4C, polynomial reduction of Y\_j(D) gives
j=+/-W mod 4L. Combining with L\textbar l gives i=+/-W mod 2L. Both i
and W are strictly between zero and L. The negative congruence cannot
hold and the positive congruence forces i=W. Therefore

\begin{verbatim}
|C|=Y_W(|A|).
\end{verbatim}

Every primary Pell index j\_i is then less than W. Write epsilon=sgn(A),
y\_i=sigma\_iY\_(j\_i)(a), and tau\_i=sigma\_i epsilon\^{}(j\_i-1) for
j\_i\textgreater0. Reduction modulo A-1 makes the pool divisible by A-1
at the finite representative E=sum(tau\_i\emph{j\_i-k\_i)\^{}2. Its
bound is
0\textless=E\textless4}15*W\^{}2\textless4W\^{}3\textless\textbar A-1\textbar,
so E=0. Thus

\begin{verbatim}
y_i=Y_(k_i)(A), x_i=+/-X_(k_i)(A)
\end{verbatim}

with signed indices, including zero. No positive-index ledger was
borrowed.

\subsubsection{Constructive completeness of the odd
helper}\label{constructive-completeness-of-the-odd-helper}

After the semantic front and frame have been fixed choose positive A as
in section 2, U=X\_W(A), C=Y\_W(A). Since W is even, U is odd and C is
even. Put H=X\_(WC)(A)=X\_C(U) and J=Y\_(WC)(A)=C\emph{Y\_C(U). The
odd-base valuation lemma above gives v2(Y\_C(U))=v2(C). Expanding the
Pell power gives Y\_C(U)=C}U\^{}(C-1)+(U\^{}2-1)T for an integer T; the
cap norm says U\^{}2-1=(A\^{}2-1)C\^{}2. Hence C divides Y\_C(U).
Therefore J/C\^{}2 is a positive odd integer and e=(J/C\^{}2-1)/2 is
integral. Also H is odd, gcd(H,J)=1, and hence gcd(H,4J)=1. CRT supplies
a sufficiently large positive D satisfying D=A mod H and D=1 mod 4J.
Choose its index-W canonical pair S,R. Polynomial base transfer gives
S=U mod H and R=W mod 4C. This supplies all comparison quotient
witnesses and the complete odd helper.

This construction reuses the old even-index odd-quotient idea, but is
now composed with the signed-target quadratic cap, where C and A may
initially have either sign and the old comparison-base lower graph is
absent.

\subsection{4. Restoring the complete accepted signed-cap
source}\label{restoring-the-complete-accepted-signed-cap-source}

It is false to restore the old helper merely by changing the sign of J:
its comparison-base lower graph does not follow for the original
negative D. Instead, after section 3 recovers
\textbar C\textbar=Y\_W(a), rebuild its private helper. Let
L=\textbar C\textbar, H=X\_(WL)(a), J=Y\_(WL)(a), and choose large
positive D by CRT with D=A mod H and D=1 mod 4J. Since W is even,
X\_W(A)=X\_W(a). Choose S to be X\_W(D) or -X\_W(D) according to the
actual sign of U, and choose R=Y\_W(D). Then S=U mod H and R=W mod 4C.
Fill all quotients and products, set the OLD e=J/C\^{}2-1, and provide
its old Dminus1-4 graph. This fixes A, C, U, every primary pair and
every semantic port, but replaces the complete private helper tuple. It
is an existential restoration, not a pointwise whole-body identity.

Restore the old frame parameter by z\_old=\textbar2z+1\textbar-1 and
reflect its two ordinate multiplication gates if necessary. Its norm and
A stay fixed, and z\_old\textgreater=0. Supply t0's old ternary graph
because t0=h0\^{}2\textgreater=0. The new support products are handled
next. Only after restoring every old graph and the full old tuple is its
whole-source theorem invoked.

\subsection{5. Acyclic sign and support
recovery}\label{acyclic-sign-and-support-recovery}

The unchanged r,s\textgreater=0 and output slack o=v185\textgreater=0
give B=r+s+1\textgreater=1 and q=B+o+1\textgreater=2. Section 3 recovers
the signed t pair without using the removed support signs. The actual
q-base request at t\textgreater=16 gives L=v181=q\^{}t and orients that
pair relative to signed A. The same pair then gives V0=(F\^{}2+1)\^{}t,
including the base-one case. The unchanged ceiling is FB=V0+s. Hence
F\textgreater=1. As in the retained fused proof, this yields
B\textgreater F\^{}2+1 and q\textgreater=3 before interpreting B as a
binomial.

The signed-A request lemma used here is elementary: with T=2An-n\^{}2-1,
H\_k=X\_k(A)+(n-A)Y\_k(A), one has H\_k=n\^{}k mod T and
H\_k\emph{(X\_k(A)-(n-A)Y\_k(A))=1 mod T. The frame implies
\textbar T\textbar\textgreater W\^{}W while both
\textbar m-n\^{}k\textbar{} and \textbar m}n\^{}k+1\textbar{} are below
W\^{}W. For n\textgreater=2,k\textgreater=1 the opposite abscissa branch
would give m\emph{n\^{}k=-1, impossible. The canonical branch gives
m=n\^{}k. For a negative target k=-l\textless0 the ordinate is -Y\_l(A).
The canonical abscissa branch is J\_l=X\_l(A)-(n-A)Y\_l(A), giving
m}n\^{}l=1 modulo T; the same strict frame bound makes this an integer
equality, impossible for n\textgreater=2. The opposite branch is -H\_l
and gives m=-n\^{}l. For k=0, y=0 and x=+/-1, so the output is m=+/-1.
These enumerate every branch. Thus a q-base output at least2 forces
k\textgreater0 and the canonical branch. They justify the unchanged x
anchor without assuming its exponent sign.

The derived t+1 recurrence now gives v189=q\^{}(t+1). The exact bound
rows166/171 say q\emph{v183+v184=L-2\textgreater0 and
q}v191+v192=v189-2\textgreater0. For q\textgreater0, uv\textgreater=0
and qu+v\textgreater0 imply u,v\textgreater=0: a negative factor would
force the other nonpositive and make the weighted sum negative. Both
support pairs are restored. Zero factors are allowed; the represented
product is uv, not uv-1.

In the cofinal-y siblings, y=V0(V0-1)\textgreater=F+1. Indeed
V0\textgreater F\^{}2+1 for t\textgreater=16,F\textgreater=1. Thus
y\textgreater=0 and r0=y-F-1\textgreater=0 without its old graph. The
q-base y+1 request gives v211=q\^{}(y+1), so row184 makes the third
weighted support sum positive and restores 213/214 from its product. The
two x support scalar guards225/226 remain unchanged: their positive
output anchor excludes nonpositive x+1 by the full signed-index lemma,
giving x\textgreater=0. No unsupported product replacement is made for
this pair.

In the conservative odd\_t sibling the unchanged scalar guards213/214
and row184 give v211=q*v213+v214+2\textgreater=2. The signed-index q
request then excludes y+1\textless=0 and gives y\textgreater=0 and
v211=q\^{}(y+1). Thus every denominator used in the eleven successor
residue ideals is also a genuine positive q-power in this sibling,
without borrowing the later r0 sign.

The unchanged successor residue ideals give every q+1-base output m=1
mod q. If its signed target k were negative, the only possible
integer-output branch above would be m=-(q+1)\^{}(-k)=-1 mod q,
impossible because q\textgreater=3. Thus all nine independent targets
v16+1,\ldots,v24+1, and the derived-label target are nonnegative before
strictness. The zero target is not excluded by this residue alone. The
product of the first eight targets is P8. In the conservative sibling
the original mixed graph forces P8\textgreater0 and r0+1\textgreater0;
in the cofinal-y siblings P8-1\textgreater=0 forces P8\textgreater0,
while r0\textgreater=0 was independently recovered above. Thus every
original scalar sign required by the complete fused source has been
restored. Its old mixed graph is supplied with a fresh triple. All other
retained witnesses are copied; none of the fresh representations is
asserted unique.

For the ninth representative, the endpoint v24=-1 would make the
unchanged raw index row B*v27=F\^{}2+1. But
0\textless F\^{}2+1\textless B: indeed
FB\textgreater=V0=(F\^{}2+1)\^{}t\textgreater F(F\^{}2+1), because
t\textgreater=16 and F\textgreater=1. No integer multiple of the
positive B lies strictly between0 and B. Hence v24+1\textgreater0 as
well. All nine representatives used in H0 are now strictly positive.

\subsection{6. Cofinal completeness, not same-tuple
equivalence}\label{cofinal-completeness-not-same-tuple-equivalence}

Begin with a finite C222 refutation and expand its proof DAG into a
finite tree. The inherited native construction permits every
sufficiently large horizon: insert truth-padding steps to reach exactly
t. Choose h0 large enough that t=16+h0\^{}2 exceeds its step count,
every local parameter and all constructor label sums. Set F=t!. This
proves completeness of the square restriction; mere unboundedness of an
unspecified set would not.

For the conservative sibling choose the ordinary beta modulus freely, as
in the parent. For cofinal-y put b=F\^{}2+1, V0=b\^{}t and y=V0(V0-1).
Every pending accumulator has at most t+1 factors, each at most
1+tF\textless=b, so it is at most b\^{}(t+1). Since t\textgreater=16,
y=b\^{}t(b\^{}t-1)\textgreater b\^{}(t+1). Also V0=1 mod F, so
F\textbar y, and y\textgreater=F+1. These are the complete inherited
beta-CRT hypotheses. A common prime of 1+(i+1)y and 1+(j+1)y would
divide i-j, hence F and y, contradiction. Thus CRT supplies x and the
endpoint quotients. Set r0=y-F-1\textgreater=0.

Construct the parent's remaining semantic data and its genuine powers in
their existing dependency order. With that front fixed, choose the odd
frame, canonical primary pairs, cap, and odd helper of sections 2--3.
For the quadratic cap slack let h=max k\_i\textless W. Then
C=Y\_W(A)\textgreater=Y\_(h+1)(A)\textgreater(2A-1)Y\_h(A). Since A is
enormous and h\textgreater=t\textgreater=16,
C\^{}2\textgreater15Y\_h(A)\^{}2+1\textgreater=sum\_i
Y\_(k\_i)(A)\^{}2+1. Thus the exact quadratic cap graph has a
nonnegative argument. Each difference Y\_(k\_i)(A)-k\_i is divisible by
A-1, so the pool quotient is integral. Supply every
retained/product/strict graph using its actual nonnegative argument.
This constructs a full zero of each source.

Together with sections 4--5, this proves soundness/completeness of the
conservative and cofinal-y siblings at the complete parent's
calibration.

\subsection{7. Optional self-sized radix
sibling}\label{optional-self-sized-radix-sibling}

This sibling has a different semantic front, not a same-front map to the
old equation q=N\^{}K. It uses the natural arithmetic argument retained
in \texttt{source/COFINAL\_131\_PROOF.md} and the same complete C222
calibration.

Let a\_i=v(15+i)+1 for 1\textless=i\textless=9. Replace the aggregate
bound by H0=1+sum a\_i and q=(B H0)\^{}K, keeping N=B\^{}K. In the
literal body v178=sum a\_i, v174=B(1+v178); the former N-square v175 now
equals v174\^{}2, and the third power base changes from v26 to v174. The
only new frame atom is v174, so the exact even frame has 38 distinct
atoms. The corresponding scalar slack graph is deleted. No base is
omitted from the frame; the old N atom remains because N is still a
requested output.

All preceding presemantic steps remain valid because they use
q\textgreater=2, not its unknown power meaning. The q/q+1 requests and
successor residue ideals establish strict positive representatives
before H0 is used as a positive base. Now H0\textgreater1 and
K=2(2+t+x+y+F)\textgreater0. The B request orients the K pair, giving
N=B\^{}K; its same-pair B*H0 request gives q=(B H0)\^{}K. Thus
N\textbar q and q\textgreater=B\^{}2\textgreater V0. The exact factorial
binomial is recovered by the inherited carry-free coefficient bound, and
the unchanged square- output fixed-point proof gives F=t! (no
consistency assumption).

Every independent a\_i is below H0. For needed lower range index
k\textless R0,
binom(a\_i,j)\textless=a\_i\^{}k\textless H0\^{}K\textless=q for
j\textless=k, with coefficient zero beyond the upper index. The derived
label has the inherited majorant Jlabel\textless=(n+1)\^{}ell,
ell=4194304, with n+1\textless=H0; padding and F=t! give
ell(t+1)\textless2F\textless K. Hence all its needed coefficients are
below q. The unchanged local Horn majorant is below
R0\^{}8388608\textless2\^{}R0\textless2\^{}K. These are precisely the
arithmetic inequalities used in the retained prime-power range decoding
and multiplicity-flow proof. No old aggregate bound or old identity
q=N\^{}K is reused. All interpreted denominators are positive, and the
retained paired intervals restore their digit signs. The natural
CRT/Horn argument therefore extracts a finite refutation.

Conversely, after the cofinal beta code fixes x,y,F and hence R0,K,B,N,
choose the nine independent positive CRT representatives at precision2K.
Only then define H0 and q=(B H0)\^{}K. The bounds just proved validate
all digits for these representatives; there is no circular future
modulus. Construct the frame and private Pell certificates last. This
supplies a complete zero and proves the same finite-refutation
equivalence.

\subsection{8. Exact compilation and mathematical verification
boundary}\label{exact-compilation-and-mathematical-verification-boundary}

The inherited Horn rewrite uses exact free-coefficient identities
product(L\_i\^{}2)=(product L\_i)\^{}2 and r\^{}2s\^{}2=(rs)\^{}2, with
complete private occurrence checks. It reconstructs old private Horn
witnesses in the reverse direction, rather than treating all private
coordinates as unchanged.

Each monic projection step solves v=f or -v=-f, where f is an integer
polynomial independent of v. It substitutes in every remaining row and
checks that all residual degrees remain at most two. Reverse the
recorded steps to restore every deleted coordinate; restore actual
unused coordinates as zero, then restore old Horn witnesses. Integer
translations old=new+shift are bijections. Finally a sum of real squares
vanishes iff every row vanishes. These give exact existential maps
through every finite compiler stage.

The independent checker replays literal edits, guard privacy, frame atom
coverage, the Horn certificate, all monic restorations, translation, and
every collected coefficient. Finite Pell, parity and support fixtures
and mutation controls supplement but do not replace the unbounded proofs
above. All branch dimensions/terms/heights are actual measured complete
equations. No remote write, repository commit, machine run or
consistency search occurs.
